% GENERATED FILE. Edit canonical lessons, then run npm run generate:books.
% canonical-source-hash: fnv1a64:ccff65d2b0646b2f
% canonical-lessons: ZH-C520-jiadian, ZH-C520-kaiguan, ZH-C520-reshuiqi, ZH-C520-nuanqi, ZH-C520-shengdian

\chapter{Around the House}
\label{ch:zh-around-the-house}
\hlchaptermodality{520}

\begin{chapteropening}
\textbf{By the end of this chapter:} \emph{I can talk about household appliances, light switches, the water heater, the heating and saving electricity.}

Complete the last lesson of chapter 520: I can talk about household appliances, light switches, the water heater, the heating and saving electricity.
\end{chapteropening}

\section[\texorpdfstring{\zh{家电} (jiādiàn)}{家电 (jiādiàn)}]{\zh{家电} (jiādiàn) — household appliance}

\label{lesson:ZH-C520-jiadian}



\begin{quote}
Before the new one: say the Mandarin for a snowy landscape, then the Mandarin for to sit in the cool.
\end{quote}

\subsection*{You'll want to know: \zh{家电}}
\textbf{\zh{家电}} — \emph{jiādiàn} — \textquotedblleft{}household appliance\textquotedblright{}.

\begin{quote}
\textbf{\zh{我们去年买了一些新家电}} \emph{wǒmen qùnián mǎi le yìxiē xīn jiādiàn} — We bought some new household appliances last year.
\end{quote}

\begin{cousinweb}
\textbf{\zh{家电}} (\emph{jiādiàn}) is built from \zh{家} 'home' and \zh{电} 'electricity', both inherited from Middle Chinese, \emph{kae} and \emph{denH} (Guangyun, in Baxter's transcription). The dictionary explains the second as the flash when yin and yang clash, lightning; the sense 'electricity' is modern.
\end{cousinweb}

\begin{grammarlens}[title={What you've built}]
One more word for saying how it was.
\end{grammarlens}

\subsection*{Your turn}
\begin{itemize}
\raggedright
  \item \emph{Say it:} \emph{jiādiàn}
  \item \emph{Say it:} \emph{jiādiàn}, once more
  \item \emph{Say it:} \emph{wǒmen qùnián mǎi le yìxiē xīn jiādiàn} — We bought some new household appliances last year.
  \item \emph{Recall:} say the Mandarin for a snowy landscape, then the Mandarin for to sit in the cool, then say \emph{jiādiàn} again
\end{itemize}

\begin{tcolorbox}[breakable,colback=teal!4,colframe=teal!35!black,title={Before you move on}]
What does \zh{家电} mean? (\textquotedblleft{}Household appliance\textquotedblright{}.) Say it once more.
\end{tcolorbox}
\section[\texorpdfstring{\zh{开关} (kāiguān)}{开关 (kāiguān)}]{\zh{开关} (kāiguān) — a light switch; a switch on a machine}

\label{lesson:ZH-C520-kaiguan}



\begin{quote}
Before the new one: say the Mandarin for to sit in the cool, then the Mandarin for a household appliance.
\end{quote}

\subsection*{You'll want to know: \zh{开关}}
\textbf{\zh{开关}} — \emph{kāiguān} — \textquotedblleft{}a light switch; a switch on a machine\textquotedblright{}.

\begin{quote}
\emph{wǒ zhǎo bu dào dēng de kāiguān} — I can't find the light switch.
\end{quote}

\begin{cousinweb}
\textbf{\zh{开关}} (\emph{kāiguān}) is built from \zh{开} 'to open' and \zh{关} 'to close', both inherited from Middle Chinese, \emph{khoj} and \emph{kwaen} (Guangyun, in Baxter's transcription), which glosses the second, after an older dictionary, as the bar laid across a door to hold it shut: the thing that opens and closes.
\end{cousinweb}

\begin{grammarlens}[title={What you've built}]
One more word for saying how it was.
\end{grammarlens}

\subsection*{Your turn}
\begin{itemize}
\raggedright
  \item \emph{Say it:} \emph{kāiguān}
  \item \emph{Say it:} \emph{kāiguān}, once more
  \item \emph{Say it:} \emph{wǒ zhǎo bu dào dēng de kāiguān} — I can't find the light switch.
  \item \emph{Recall:} say the Mandarin for to sit in the cool, then the Mandarin for a household appliance, then say \emph{kāiguān} again
\end{itemize}

\begin{tcolorbox}[breakable,colback=teal!4,colframe=teal!35!black,title={Before you move on}]
What does \zh{开关} mean? (\textquotedblleft{}A light switch; a switch on a machine\textquotedblright{}.) Say it once more.
\end{tcolorbox}
\section[rèshuǐqì]{rèshuǐqì — water heater}

\label{lesson:ZH-C520-reshuiqi}



\begin{quote}
Before the new one: say the Mandarin for a household appliance, then the Mandarin for a light switch.
\end{quote}

\subsection*{You'll want to know: rèshuǐqì}
\emph{rèshuǐqì} — \textquotedblleft{}water heater\textquotedblright{}.

\begin{quote}
\emph{wǒmen de rèshuǐqì yòng le shí nián le} — Our water heater has been going for ten years.
\end{quote}

Its characters come later; for now it is a word for the ear.

\begin{cousinweb}
\emph{rèshuǐqì} is built from \zh{热} 'hot', \zh{水} 'water' and \emph{qì} 'a vessel', all inherited from Middle Chinese, \emph{nyet}, \emph{sywijX} and \emph{khijH} (Guangyun, in Baxter's transcription), which glosses the last 'vessels and dishes'.
\end{cousinweb}

\begin{grammarlens}[title={What you've built}]
One more word for saying how it was.
\end{grammarlens}

\subsection*{Your turn}
\begin{itemize}
\raggedright
  \item \emph{Say it:} \emph{rèshuǐqì}
  \item \emph{Say it:} \emph{rèshuǐqì}, once more
  \item \emph{Say it:} \emph{wǒmen de rèshuǐqì yòng le shí nián le} — Our water heater has been going for ten years.
  \item \emph{Recall:} say the Mandarin for a household appliance, then the Mandarin for a light switch, then say \emph{rèshuǐqì} again
\end{itemize}

\begin{tcolorbox}[breakable,colback=teal!4,colframe=teal!35!black,title={Before you move on}]
What does rèshuǐqì mean? (\textquotedblleft{}Water heater\textquotedblright{}.) Say it once more.
\end{tcolorbox}
\section[nuǎnqì]{nuǎnqì — central heating}

\label{lesson:ZH-C520-nuanqi}



\begin{quote}
Before the new one: say the Mandarin for a light switch, then the Mandarin for a water heater.
\end{quote}

\subsection*{You'll want to know: nuǎnqì}
\emph{nuǎnqì} — \textquotedblleft{}central heating\textquotedblright{}.

\begin{quote}
\emph{yǒu nuǎnqì de fángjiān bǐ wàibian nuǎnhuo} — A room with heating is warmer than outside.
\end{quote}

Its characters come later; for now it is a word for the ear.

\begin{cousinweb}
\emph{nuǎnqì} is built from \emph{nuǎn} 'warm' and \zh{气} 'air', both inherited from Middle Chinese, \emph{nwanX} and \emph{khj+jH} (Guangyun, in Baxter's transcription). The dictionary gives the first as the same word as another spelling glossed 'the warmth of a fire', and the second 'breath'.
\end{cousinweb}

\begin{grammarlens}[title={What you've built}]
One more word for saying how it was.
\end{grammarlens}

\subsection*{Your turn}
\begin{itemize}
\raggedright
  \item \emph{Say it:} \emph{nuǎnqì}
  \item \emph{Say it:} \emph{nuǎnqì}, once more
  \item \emph{Say it:} \emph{yǒu nuǎnqì de fángjiān bǐ wàibian nuǎnhuo} — A room with heating is warmer than outside.
  \item \emph{Recall:} say the Mandarin for a light switch, then the Mandarin for a water heater, then say \emph{nuǎnqì} again
\end{itemize}

\begin{tcolorbox}[breakable,colback=teal!4,colframe=teal!35!black,title={Before you move on}]
What does nuǎnqì mean? (\textquotedblleft{}Central heating\textquotedblright{}.) Say it once more.
\end{tcolorbox}
\section[shěngdiàn]{shěngdiàn — to use little electricity; to save electricity}

\label{lesson:ZH-C520-shengdian}



\begin{quote}
Before the new one: say the Mandarin for a water heater, then the Mandarin for central heating.
\end{quote}

\subsection*{You'll want to know: shěngdiàn}
\emph{shěngdiàn} — \textquotedblleft{}to use little electricity; to save electricity\textquotedblright{}.

\begin{quote}
\emph{xīn jiādiàn bǐ jiù de shěngdiàn} — New appliances use less electricity than old ones.
\end{quote}

Its characters come later; for now it is a word for the ear.

\begin{cousinweb}
\emph{shěngdiàn} is built from \emph{shěng} 'to spare' and \zh{电} 'electricity', both inherited from Middle Chinese, \emph{srjaengX} and \emph{denH} (Guangyun, in Baxter's transcription). The dictionary glosses the first as a government office and, under another reading, 'to examine'; the sense 'to spare, to save' is not given there.
\end{cousinweb}

\begin{grammarlens}[title={What you've built}]
One more word for saying how it was.
\end{grammarlens}

\subsection*{Your turn}
\begin{itemize}
\raggedright
  \item \emph{Say it:} \emph{shěngdiàn}
  \item \emph{Say it:} \emph{shěngdiàn}, once more
  \item \emph{Say it:} \emph{xīn jiādiàn bǐ jiù de shěngdiàn} — New appliances use less electricity than old ones.
  \item \emph{Recall:} say the Mandarin for a water heater, then the Mandarin for central heating, then say \emph{shěngdiàn} again
\end{itemize}

\begin{tcolorbox}[breakable,colback=teal!4,colframe=teal!35!black,title={Before you move on}]
What does shěngdiàn mean? (\textquotedblleft{}To use little electricity; to save electricity\textquotedblright{}.) Say it once more.
\end{tcolorbox}
